\documentclass{article}
\usepackage[UTF8]{ctex}
\usepackage{amsmath,amssymb,amsfonts}
\usepackage{enumitem}
\usepackage{geometry}
\geometry{a4paper, margin=1in}
% === 错题本排版优化（不参与题目标识） ===
\setlength{\parindent}{0pt}
\setlength{\parskip}{0.35em plus 0.10em minus 0.08em}
\setlength{\emergencystretch}{3em}
\setlist{itemsep=0.25em, topsep=0.40em, parsep=0pt}
\allowdisplaybreaks[2]
% === 错题本排版优化结束 ===
\begin{document}
\section*{上海大学2024年高等代数真题}
\section*{38．上海大学2024年高等代数试卷}
一．填空题．\\

\subsection{题目1}
1．设 \(\alpha\) 为 \(n\) 维单位实列向量，\(I\) 为单位矩阵，则 \(\left|2 I-\alpha \alpha^{\mathrm{T}}\right|=\) \(\_\_\_\_\)．

\subsection{题目2}
2．设非零矩阵 \(A\) 的秩 \(r(A)=p\) ，且 \(A^{2}=A\) ，则 \(r(A-I)+r(A+I)=\) \(\_\_\_\_\)．

\subsection{题目3}
3．设向量组 \(\alpha_{1}, \alpha_{2}, \alpha_{3}\) 线性无关，且

\[
\beta_{1}=\alpha_{1}+2 \alpha_{2}+3 \alpha_{3}, \beta_{2}=2 \alpha_{1}+2 \alpha_{2}+3 \alpha_{3}, \beta_{3}=3 \alpha_{1}+2 \alpha_{2}+3 \alpha_{3} .
\]

则 \(\beta_{1}, \beta_{2}, \beta_{3}\) 的极大线性无关组为 \(\_\_\_\_\)．

\subsection{题目4}
4．设 \(A=\left(\begin{array}{lll}1 & 3 & a \\ 2 & 2 & b \\ 0 & 0 & 5\end{array}\right)\) ，其中 \(a, b\) 为常数．设线性空间 \(V=\left\{B \in \mathbb{C}^{3 \times 3} \mid A B=B A\right\}\) ，则 \(V\) 的维数为\\
\(\_\_\_\_\)．

\subsection{题目5}
5．设 \(A, C\) 为 \(n\) 阶实可逆矩阵，且 \(A\) 为对称矩阵，则实对称矩阵 \(\left(\begin{array}{cc}A & C \\ C^{\mathrm{T}} & O\end{array}\right)\) 的正负惯性指数差为\\
\(\_\_\_\_\)．

二．是非题．若错误，请举出反例；若正确，请给出相应理由．

\begin{itemize}

\subsection{题目6}
  \item[6．] 多项式 \(x^{4}+x^{3}+x^{2}+x+1\) 在有理数域上不可约。

\subsection{题目7}
  \item[7．] 设 \(A\) 是 \(n\) 阶矩阵，齐次线性方程组 \(A X=0\) 有唯一解的充分必要条件是 \(|A|=0\) 。

\subsection{题目8}
  \item[8．] 设 \(A, B, C, D\) 为 \(n\) 阶矩阵，则 \(\left|\begin{array}{cc}A & B \\ C & D\end{array}\right|=|A D-C B|\) 。

\subsection{题目9}
  \item[9．] 欧氏空间上对称变换在任意一组基下的矩阵为实对称矩阵。

\subsection{题目10}
  \item[10．] 设四阶矩阵
\[
A=\left(\begin{array}{cccc}
1 & 1 & 0 & 0 \\
-1 & -1 & 0 & 0 \\
-2 & -2 & 2 & 1 \\
1 & 1 & -1 & 0
\end{array}\right)
\]
则 \(A\) 的初等因子为 \(\lambda, \lambda,(\lambda-1)^{2}\) ．
\end{itemize}

三。解答题。\\

\subsection{题目11}
11．设 \(A=\left(\begin{array}{llll}3 & 0 & 0 & 0 \\ 2 & 3 & 0 & 0 \\ 2 & 2 & 3 & 0 \\ 2 & 2 & 2 & 3\end{array}\right)\) ，且 \(A B=B+I\) ，求矩阵 \(B^{n}\) ，其中 \(n\) 为整数．

\subsection{题目12}
12．设 \(n(n>1)\) 阶矩阵

\[
A=\left(\begin{array}{ccccc}
0 & 1 & 1 & \cdots & 1 \\
1 & 0 & 1 & \cdots & 1 \\
1 & 1 & 0 & \cdots & 1 \\
\vdots & \vdots & \vdots & & \vdots \\
1 & 1 & 1 & \cdots & 0
\end{array}\right)
\]

求 \(A^{-1}\) ．

\subsection{题目13}
13．求解线性方程组

\[
\left\{\begin{array}{l}
x_{1}+2 x_{2}+3 x_{3}+4 x_{4}=1 \\
x_{1}+3 x_{2}-x_{3}-x_{4}=0 \\
3 x_{1}+8 x_{2}+\left(a^{2}-5 a+7\right) x_{3}+a x_{4}=1
\end{array}\right.
\]

其中 \(a\) 为常数．

\subsection{题目14}
14．设 \(a\) 为整数，且二次型 \(f=2 x_{1}^{2}+2 x_{2}^{2}+a x_{3}^{2}+2 x_{1} x_{2}+2 x_{1} x_{3}+2 x_{3} x_{3}\) 有三个不同整数特征值，用正交变换将此二次型化为标准形（需要写出正交变换及标准形）。

\subsection{题目15}
15．设 \(A\) 为数域 \(F\) 上的 \(n\) 阶矩阵，\(p(x)\) 为 \(F\) 上首1的不可约多项式。若存在正整数 \(r\) 使得 \(A\) 在 \(F\)上最小多项式为 \(p^{r}(x)\) ，求矩阵 \(P=\left(\begin{array}{cc}A & I \\ O & A\end{array}\right)\) 的最小多项式。

\subsection{题目16}
16．设 \(V\) 是有限维复空间，\(S, T\) 是 \(V\) 上的两个线性变换，满足 \(S T-T S=S^{2}-S\) 。对于任意复数 \(\mu\)，定义 \(V\) 的子空间 \(V_{\mu}\) 如下：

\[
V_{\mu}=\left\{\alpha \in V \mid \text { 存在正整数 } m \text { 使得 }(S-\mu I)^{m} \alpha=0\right\} \text {. }
\]

其中 \(I\) 为恒等变换．

\begin{itemize}
  \item[（1）] 求证 \(V_{\mu}\) 是 \(T\) 的不变子空间．
  \item[（2）] 如果 \(V_{\mu} \neq\{0\}\) ，求证 \(\mu=0\) 或者 1 。
\end{itemize}
\end{document}
